\documentclass[10pt,a4paper]{article}
\usepackage[utf8]{inputenc}
\usepackage[a4paper,left=3cm,right=2cm]{geometry}
\usepackage{amsmath}
\usepackage[thmmarks, amsmath, thref]{ntheorem}
\usepackage{amsfonts}
\usepackage{amssymb}
\usepackage{fancyhdr}
\pagestyle{fancy}
\usepackage{anysize}
\usepackage{graphicx}
\usepackage{float}
\usepackage{hyperref}
\usepackage{multirow}
\usepackage{enumitem}
\newtheorem{theorem}{Theorem}
\theoremsymbol{\ensuremath{\blacksquare}}
\newtheorem*{solution}{Solución:}
\usepackage{listings}
\usepackage{xcolor}

% Margenes y formalidades

\textwidth 6.25in % Margen Derecho y ancho del texto
\textheight 8.5in
\oddsidemargin 0in % Margen Izquierdo
\headheight 0in % Largo del texto
\leftmargin 1in
\parindent 0em % Salto entre la indentación (Sangría)https://www.overleaf.com/4519127669sdpptkhvsbfn
\topmargin -0.5in % Margen Superior
\parskip 0.5ex % Salto entre parrafos
\baselineskip 1pt

%\lhead{\includegraphics[scale=0.2]{logo.png}} % texto izquierda de la cabecera
%\chead{}                                      % texto centro de la cabecera
%\rhead{\includegraphics[scale=0.2]{FIG/logo-mba-2016.png}} % número de página a la derecha

\renewcommand{\headrulewidth}{0.4pt} % grosor de la línea de la cabecera
\renewcommand{\footrulewidth}{0.4pt} % grosor de la línea del pie

\begin{document}



\pagebreak
\vspace*{4mm} 

\begin{center}
\begin{huge}
\textbf{\\Econometría: Ejercicios 1}
\end{Large}\\
%{\large 2023-2}\\
\end{center}



\section{Estadística Descriptiva}

\textbf{Ejercicio 1:} Se tiene el siguiente conjunto de datos sobre el número de ventas realizadas por 5 vendedores en un mes:
$$X = \{10,15,12,18,20\}$$
Calcule la media, mediana y varianza muestral.

{\color{red} \textbf{Solución:}

La media es:
$$\Bar{X} = \frac{10+15+12+18+20}{5} = \frac{75}{5} = 15$$

Ordenando los datos: {10, 12, 15, 18, 20}, la mediana es el valor central, es decir:

$$mediana(X) = 15$$

La varianza muestral se calcula como:

$$s^2 = \frac{\sum (X_i - \Bar{X})^2}{n-1}$$

$$s^2 = \frac{(10-15)^2 + (15-15)^2 + (12 - 15)^2 + (18-15)^2 + (20-15)^2}{4}$$

$$s^2 = \frac{(-5)^2+(0)^2 + (-3)^2 +(3)^2 + (5)^2}{4} = \frac{25+0+9+9+25}{4} = \frac{68}{4} = 17$$}








\section{Muestreo}

% \begin{itemize}
%   \item \textbf{Muestreo aleatorio simple (MAS):} todas las unidades tienen la misma probabilidad de selección.
%     \item \textbf{Sistemático:} se ordena la población y se elige un inicio aleatorio; luego cada $k$-ésima unidad. Es eficiente si no hay periodicidad alineada con $k$.
%   \item \textbf{Estratificado:} la población se divide en estratos homogéneos; se muestrea aleatoriamente dentro de cada estrato y se pondera por tamaños.
%   \item \textbf{Conglomerados:} se seleccionan grupos (heterogéneos internamente) y se censa o muestrea dentro de ellos; útil cuando listar individuos es costoso.
% \end{itemize}

\textbf{Ejercicio 2:} Se tiene una población de 1000 individuos con una media de ingreso mensual de $1500$ y una desviación estándar de $300$. Si se toma una muestra aleatoria simple de 100 individuos, ¿cuál es la media esperada y la desviación estándar de la media muestral?

{\color{red} \textbf{Solución:}

La media esperada de la muestra es igual a la media poblacional:
$$E(\Bar{X}) = \mu = 1500$$


La desviación estándar de la media muestral (error estándar) se calcula como:

$$\sigma_{\Bar{X}} = \frac{\sigma}{\sqrt{n}} = \frac{300}{\sqrt{100}} = \frac{300}{10} = 30$$}



\section{Intervalos de confianza}



\textbf{Ejercicio 3:} Josefina se encuentra trabajando para una importante constructora nacional en el \'area de finanzas, el jefe de
proveedores le comenta que la media de cuentas por cobrar es de 220.000.000 de pesos. Josefina tiene una  muestra aleatoria de 500 facturas con media igual a 195.000.000 de pesos y una desviaci\'on
est\'andar de 1.000.000. Con un 90\% de confianza ¿es correcta la afirmacion del jefe?


    $$\color{red}{\bar{x} \pm z_{\frac{\alpha}{2}}\frac{\sigma}{\sqrt{n}}}$$

{\color{red}$[194.926.209,8 ; 195.073.790,2]$ }


   {\color{red}  El valor que comenta el jefe es 220.000.000 el cual no está en el intervalo de confianza calculado, así que  90\% de confianza la afirmación es incorrecta. }
\\
   
\textbf{Ejercicio 4:} Una empresa desea comparar los tiempos promedio de producción de dos plantas, Planta A y Planta B. Se toman muestras aleatorias de los tiempos de producción (en horas) de ambas plantas.

\begin{itemize}
    \item Para la \textbf{Planta A}, se selecciona una muestra de 50 observaciones con una media muestral de \( \bar{X}_A = 5,2 \) horas y una varianza conocida de \( \sigma_A^2 = 0,16 \) horas.
    \item Para la \textbf{Planta B}, se selecciona una muestra de 60 observaciones con una media muestral de \( \bar{X}_B = 4,9 \) horas y una varianza conocida de \( \sigma_B^2 = 0,25 \) horas.
\end{itemize}

 Verifique si la diferencia entre los tiempos de producción de las dos plantas es estadísticamente significativa.

 {\color{red}
 El intervalo de confianza para la diferencia de medias con varianzas conocidas se calcula con la siguiente fórmula:

 \[
 (\bar{X}_A - \bar{X}_B) \pm Z_{\alpha/2} \cdot \sqrt{\frac{\sigma_A^2}{n_A} + \frac{\sigma_B^2}{n_B}}
\]

{\bf Paso 1: Diferencia de medias}

 \[
 \bar{X}_A - \bar{X}_B = 5.2 - 4.9 = 0.3
 \]

{\bf Paso 2: Error estándar de la diferencia de medias}

 \[
 \text{SE} = \sqrt{\frac{0.16}{50} + \frac{0.25}{60}} = \sqrt{0.0032 + 0.004167} = \sqrt{0.007367} \approx 0.0858
 \]

{\bf Paso 3: Cálculo del intervalo de confianza}

\[
 0.3 \pm 2.576 \times 0.0858
 \]

 \[
 0.3 \pm 0.2209
 \]


 El intervalo de confianza al 99\% para la diferencia de medias es:

\[
(0.0791, 0.5209)
\]


Con un 99\% de confianza, la diferencia en los tiempos promedio de producción entre las dos plantas está entre \(0.0791\) y \(0.5209\) horas. Como el valor 0 no está dentro del intervalo, se puede concluir que la diferencia es estadísticamente significativa al nivel de confianza del 99\%.
}


\section{Regresión Lineal Simple}
\textbf{Ejercicio 5:} Se tienen los siguientes datos de salario, años de educación y experiencia laboral de 5 individuos:
\begin{center}
\begin{tabular}{|c|c|c|c|}
\hline
\textbf{Individuo} & \textbf{Educ} ($X_1$) & \textbf{Exp} ($X_2$)  & \textbf{Salario ($Y$)}\\
\hline
1 & 12 & 14 & 1000 \\
2 & 14 & 15 & 1200 \\
3 & 15 & 17 & 1300 \\
4 & 16 & 15 & 1200 \\
5 & 18 & 16 & 1500 \\
\hline
\end{tabular}
\end{center}

\begin{enumerate}
    \item Planteando un modelo donde el salario sólo dependa de la educación se tiene la siguiente ecuación:
    $$\widehat{Salario} = \beta_0 + \beta_1 \cdot Educ$$
    Calcule los valores de los estimadores $\hat{\beta_0}$ y $\hat{\beta_1}$ y SST, SSE, SSR, $R^2$.
    \\
    
    \textbf{solución}


\textcolor{red}{

\[
\bar X=\frac{12+14+15+16+18}{5}=15,\qquad
\bar Y=\frac{1000+1200+1300+1200+1500}{5}=1240.
\]

\begin{center}
\begin{tabular}{|r|r|r|r|r|r|r|}
\hline
$X_i$ & $Y_i$ & $X_i-\bar X$ & $Y_i-\bar Y$ & $(X_i-\bar X)^2$ & $(Y_i-\bar Y)^2$ & $(X_i-\bar X)(Y_i-\bar Y)$\\ \hline
12 & 1000 & $-3$ & $-240$ & $9$  & $57{,}600$ & $720$ \\
14 & 1200 & $-1$ & $-40$  & $1$  & $1{,}600$  & $40$  \\
15 & 1300 & $0$  & $60$   & $0$  & $3{,}600$  & $0$   \\
16 & 1200 & $1$  & $-40$  & $1$  & $1{,}600$  & $-40$ \\
18 & 1500 & $3$  & $260$  & $9$  & $67{,}600$ & $780$ \\ \hline
\multicolumn{2}{|r|}{\textbf{Totales}} & & & $\mathbf{20}$ & $\mathbf{132{,}000}$ & $\mathbf{1{,}500}$\\ \hline
\end{tabular}
\end{center}

\[
S_{xx}=20,\quad S_{xy}=1{,}500.
\]


\textbf{Estimadores MCO:}
\[
\hat\beta_1=\frac{S_{xy}}{S_{xx}}=\frac{1500}{20}=75,\qquad
\hat\beta_0=\bar Y-\hat\beta_1\bar X=1240-75\cdot 15=115.
\]
\[
\widehat{Y}_i=115+75\,X_i.
\]

\textbf{Descomposición de suma de cuadrados y } 

\begin{center}
\begin{tabular}{|r|r|r|r|r|}
\hline
$Educ$ & $Y$ & $\widehat{Y}=115+75\,Educ$ & $e=Y-\widehat{Y}$ & $e^{2}$ \\ \hline
12 & 1000 & 1015 & $-15$  & $225$ \\
14 & 1200 & 1165 & $35$   & $1225$ \\
15 & 1300 & 1240 & $60$   & $3600$ \\
16 & 1200 & 1315 & $-115$ & $13225$ \\
18 & 1500 & 1465 & $35$   & $1225$ \\ \hline
\multicolumn{4}{|r|}{\textbf{SSR} $=\sum e_i^2$} & \textbf{19500} \\ \hline
\end{tabular}
\end{center}



$R^2$:
\[
SST=132{,}000,\qquad SSR=19{,}500
\]
\[
SSE=SST-SSR=132{,}000-19{,}500=12{,}500,\qquad
R^2=\frac{SSR}{SST}=\frac{112{,}500}{132{,}000}\approx 0.8523.
\]

\textbf{Lectura:} $\widehat{Salario}=115+75\cdot Educ$. Un año adicional de educación se asocia, en promedio, con \$75 más de salario. El ajuste explica $\approx 85.23\%$ de la variación de $Y$.}



\end{enumerate}


\section{Formulario}
\begin{itemize}

\item El intervalo de confianza de $100(1-\alpha)\%$ para la media poblacional $\mu$ cuando $\sigma$ es \emph{\textbf{conocido}}:

$$\color{red}{\bar{x} \pm z_{\frac{\alpha}{2}}\frac{\sigma}{\sqrt{n}}}$$

\item El intervalo de confianza de $100(1-\alpha)\%$ para la diferencia de medias
$\mu_x-\mu_y$ es:

$${\color{red}(\bar{X}-\bar{Y}) \pm z_{\frac{\alpha}{2}}\sqrt{\frac{\sigma_x^2}{n_x}+\frac{\sigma_y^2}{n_y}}}$$


\item El intervalo de confianza de $100(1-\alpha)\%$ para la media poblacional $\mu$ cuando $\sigma$ es \emph{\textbf{desconocido}}:

$$\color{red}{\bar{x} \pm t_{(n-1)\frac{\alpha}{2}}\frac{S}{\sqrt{n}}}$$

\item  Estimadores MCO para RLS.
{
$$\hat{\beta}_0 = \bar{y} - \hat{\beta}_1 \bar{x} $$
$$\hat{\beta}_1 = \frac{\sum_{i=1}^{n}(x_i - \bar{x})(y_i-\bar{y})}{\sum_{i=1}^{n}(x_i - \bar{x})^2}=\frac{Cov(x,y)}{Var(x)}$$}

\item \[ R^2  = 1 - \frac{SSR}{SST} \] 
\item  $$SST = \sum_{i=1}^{n} (y_i - \bar{y})^2 \quad \quad SSR = \sum_{i=1}^{n}  \hat{u}_i^2 $$
    \item \[ \bar{R}^2 =  1 - \frac{(n-1)SSR}{(n-k-1)SST} \]

\end{itemize}

\begin{figure}
    \centering
    \includegraphics[width=0.8\linewidth]{Figuras/normal_critic.png}
\end{figure}





\end{document}